\chapter*{Tài liệu tham khảo}
\addcontentsline{toc}{section}{Tài liệu tham khảo}

\begin{enumerate}
\item Akenine-Möller, T., Haines, E., \& Hoffman, N. (2018). \textit{Real-Time Rendering} (4th ed.). A K Peters/CRC Press.

\item Bộ Giáo dục và Đào tạo. (2018). \textit{Chương trình giáo dục phổ thông môn Vật lý}. Ban hành kèm Thông tư 32/2018/TT-BGDĐT.

\item Bộ Giáo dục và Đào tạo Việt Nam. (2018). \textit{Chương trình giáo dục phổ thông môn Vật lý}. Nhà xuất bản Giáo dục Việt Nam.

\item Bommasani, R., et al. (2021). On the opportunities and risks of foundation models. \textit{arXiv preprint arXiv:2108.07258}.

\item Brown, T. B., et al. (2020). Language Models are Few-Shot Learners. \textit{Advances in Neural Information Processing Systems}, 33, 1877–1901.

\item Bybee, R. W. (2013). \textit{The case for STEM education: Challenges and opportunities}. NSTA Press.

\item Chiu, T. K. F., et al. (2023). Systematic literature review on opportunities, challenges, and future research recommendations of artificial intelligence in education. \textit{Computers and Education: Artificial Intelligence}, 4, 100118.

\item de Jong, T., \& van Joolingen, W. R. (1998). Scientific discovery learning with computer simulations of conceptual domains. \textit{Review of Educational Research}, 68(2), 179–201.

\item Dirksen, J. (2023). \textit{Learn Three.js: Program 3D animations and visualizations for the web with JavaScript and TypeScript} (4th ed.). Packt Publishing.

\item Dori, Y. J., \& Belcher, J. (2005). How does technology-enabled active learning affect undergraduate students' understanding of electromagnetism concepts? \textit{The Journal of the Learning Sciences}, 14(2), 243–279.

\item Duit, R., \& Treagust, D. F. (2003). Conceptual change: A powerful framework for improving science teaching and learning. \textit{International Journal of Science Education}, 25(6), 671–688.

\item Foley, van Dam, Feiner, \& Hughes. (1996). \textit{Computer Graphics: Principles and Practice}. Addison-Wesley.

\item Galán, D., \& Herrera, J. (2016). Performance analysis of HTML5 Canvas and SVG for web-based animation and visualization. \textit{Journal of Web Engineering}, 15(3–4), 231–254.

\item Haines, E., \& Akenine-Möller, T. (2018). \textit{Real-Time Rendering} (4th ed.). CRC Press.

\item Hearn, D., \& Baker, M. P. (2011). \textit{Computer Graphics with OpenGL} (4th ed.). Pearson.

\item Holmes, W., Bialik, M., \& Fadel, C. (2023). \textit{Artificial intelligence in education: Promise and implications for teaching and learning}. Center for Curriculum Redesign.

\item Hughes, J. F., et al. (2014). \textit{Computer Graphics: Principles and Practice} (3rd ed.). Addison-Wesley.

\item Khronos Group. \textit{WebGL Specification}. Truy cập từ https://www.khronos.org/webgl/

\item Linn, M. C., \& Eylon, B. S. (2011). \textit{Science learning and instruction: Taking advantage of technology to promote knowledge integration}. Routledge.

\item Liu, P., et al. (2023). Pre-train, Prompt, and Predict: A Systematic Survey of Prompting Methods in Natural Language Processing. \textit{ACM Computing Surveys}, 55(9), 1–35.

\item Mayer, R. E. (2009). \textit{Multimedia Learning} (2nd ed.). Cambridge University Press.

\item Mozilla Developer Network. (2023). \textit{HTML5 Canvas: 2D graphics and animation}. Truy cập từ https://developer.mozilla.org/en-US/docs/Web/API/Canvas\_API

\item National Research Council. (2012). \textit{A framework for K-12 science education: Practices, crosscutting concepts, and core ideas}. National Academies Press.

\item Parisi, T. (2012). \textit{WebGL: Up and Running}. O'Reilly Media.

\item Planinic, M., Ivanjek, L., \& Susac, A. (2013). Comparison of university students' understanding of graphs in different contexts. \textit{Physical Review Special Topics - Physics Education Research}, 9(2), 020103.

\item Rutten, N., van Joolingen, W. R., \& van der Veen, J. T. (2012). The learning effects of computer simulations in science education. \textit{Computers \& Education}, 58(1), 136–153.

\item Sommerville, I. (2015). \textit{Software Engineering} (10th ed.). Pearson.

\item Sweller, J. (2011). Cognitive load theory. In J. P. Mestre \& B. H. Ross (Eds.), \textit{The psychology of learning and motivation: Cognition in education} (pp. 37–76). Elsevier Academic Press.

\item Sweller, J., van Merriënboer, J. J., \& Paas, F. (2019). Cognitive architecture and instructional design: 20 years later. \textit{Educational Psychology Review}, 31(2), 261–292.

\item Team, G., et al. (2023). Gemini: A Family of Highly Capable Multimodal Models. \textit{arXiv preprint arXiv:2312.11805}.

\item Vygotsky, L. S. (1978). \textit{Mind in society: The development of higher psychological processes}. Harvard University Press.

\item Wieman, C. E., \& Perkins, K. K. (2005). Transforming physics education. \textit{Physics Today}, 58(11), 36–41.

\item Zhang, K., \& Aslan, A. B. (2021). AI technologies for education: Recent research \& future directions. \textit{Computers and Education: Artificial Intelligence}, 2, 100025.

\end{enumerate}

